For the amplitude, we make the functional-family derivative explicit.
Suppress the endpoint type and write
\[
 H_t^0=H_t(0),\qquad G_t=(H_t^0)^{-1},\qquad
 \Delta H_t(u)=H_t(u)-H_t^0,\qquad
 \mathcal T_t(u)=G_t\Delta H_t(u).
\]
The uniformly positive tridiagonal reference operator is $C^1$ in
operator norm.  The resolvent identity gives
\begin{equation}\label{eq:v49-moving-reference}
 \dot G=-G\dot H^0G,\qquad
 \dot{\Delta H}=\dot H(u)-\dot H^0,\qquad
 \dot{\mathcal T}=\dot G\,\Delta H+G\dot{\Delta H}.
\end{equation}
In particular the derivative of the moving reference is not omitted.
For each fixed endpoint order $m$, locality and the weighted half-line
estimates imply, on a uniform small neighborhood,
\[
 \sum_{i,j\ge1}\bigl|
       \partial_u^r\partial_t^a(\Delta H_t(u))_{ij}\bigr|
       \le C_m\sum_{i\ge1}\rho^i<\infty,
          \qquad 0\le r\le m,\quad a=0,1.
\]
Only neighboring entries occur.  The terms without endpoint decay in
$\dot H(u)$ cancel those in $\dot H^0$; the remaining terms contain
an orbit coordinate or its parameter derivative.  Fixed index powers
are absorbed in the strict exponential margin.  Finite entry sums
therefore converge with their derivatives uniformly in trace norm.
Thus $\mathcal T_t$ is $C^1$ in trace norm, in each fixed endpoint norm,
and $\|\mathcal T_t\|\le q<1$ after the collar is shrunk once.

The edge sum $U_t$ of \eqref{eq:v4-half-product} has the same
summability: $g_t[-\ell_{t,uv}(0,0)]=1$ identically in $t$, so its
reference derivative cancels before the edges are summed.  In the
absolutely convergent trace series \eqref{eq:v4-half-amplitude},
cyclicity and the trace-ideal inequality give
\[
 \frac{d}{dt}\operatorname{tr}(\mathcal T_t^n)
    =n\operatorname{tr}(\mathcal T_t^{n-1}\dot{\mathcal T}_t),
 \qquad
 |\operatorname{tr}(\mathcal T_t^{n-1}\dot{\mathcal T}_t)|
     \le q^{n-1}\|\dot{\mathcal T}_t\|_1.
\]
Summing this geometric majorant proves
\begin{equation}\label{eq:v49-amplitude-derivative}
 \frac{\dot B}{B}=\dot U-
       \operatorname{tr}\bigl((I+\mathcal T)^{-1}
                                      \dot{\mathcal T}\bigr).
\end{equation}
Fixed endpoint differentiation introduces only fixed polynomial factors
in the series index and bounded resolvent factors, which remain
summable.  This proves the amplitude's $C^1$ dependence at each finite
smooth order, directly from the same relative construction.  The
half-line actions have the $C^1$ dependence already established above.
